\chapter{Cadre général du projet}

\section*{Introduction}
Ce premier chapitre pose le cadre dans lequel s'inscrit notre projet de fin d'études.
Nous présentons successivement l'entreprise d'accueil, la problématique identifiée,
l'étude critique des solutions existantes, puis la solution proposée et la méthodologie
de développement retenue.

% ============================================================
\section{Présentation de l'entreprise TAC-TIC}

\subsection{Identité et activité}

\textbf{TAC-TIC}, filiale du Groupe New Ways, est une entreprise de services
numériques (ESN) fondée en 2012 et implantée au Technoparc El Ghazela, Ariana
(Tunisie). Elle compte plus d'une centaine de collaborateurs et est certifiée
\textbf{ISO 27001} (sécurité de l'information) et \textbf{ISO 9001} (management
de la qualité). Son cœur de métier couvre le développement web et mobile,
l'architecture cloud et l'intégration de solutions d'intelligence artificielle
et d'IoT.

\subsection{Évolution historique}

\begin{table}[H]
  \centering
  \caption{Jalons historiques de TAC-TIC}
  \label{tab:historique_tactic}
  \begin{tabularx}{\textwidth}{C{2cm}X}
    \toprule
    \thead{Année} & \thead{Étape clé} \\
    \midrule
    \rowcolor{lightgray}
    2012 & Création de TAC-TIC au sein du Groupe New Ways ; premiers projets web \\
    2014 & Obtention de la certification ISO 27001 ; spécialisation en cybersécurité \\
    \rowcolor{lightgray}
    2017 & Transition vers l'architecture cloud (AWS, Azure) \\
    2020 & Lancement de la division IA/IoT ; déploiement de premiers modèles ML en production \\
    \rowcolor{lightgray}
    2023 & Dépassement du cap des 100 projets livrés ; ouverture d'un bureau à Paris \\
    2026 & Développement de TAC-TIC IntelliPM (présent PFE) \\
    \bottomrule
  \end{tabularx}
\end{table}

\subsection{Domaines d'activité}

\begin{table}[H]
  \centering
  \caption{Domaines d'activité de TAC-TIC}
  \label{tab:domaines_tactic}
  \begin{tabularx}{\textwidth}{L{3.5cm}X}
    \toprule
    \thead{Domaine} & \thead{Activités principales} \\
    \midrule
    \rowcolor{lightgray}
    Développement web & Applications SPA React/Vue, APIs REST Laravel/Node.js \\
    Développement mobile & Applications iOS/Android (React Native, Flutter) \\
    \rowcolor{lightgray}
    Cloud \& DevOps & Architecture AWS/Azure, CI/CD, Kubernetes, monitoring \\
    Intelligence artificielle & ML supervisé/non supervisé, NLP, vision par ordinateur \\
    \rowcolor{lightgray}
    Conseil \& audit & Audit sécurité ISO 27001, conseil en transformation numérique \\
    \bottomrule
  \end{tabularx}
\end{table}

\subsection{Organigramme}

\begin{figure}[H]
  \centering
  \IfFileExists{figures/organigramme_tactic.png}{%
    \includegraphics[width=0.85\textwidth]{figures/organigramme_tactic.png}%
  }{%
    \fbox{\parbox{0.85\textwidth}{\centering\vspace{1.5cm}
      \textbf{[IMAGE A INSÉRER]}\\[0.4em]
      \texttt{figures/organigramme\_tactic.png}\\[0.3em]
      \textit{Organigramme de l'entreprise TAC-TIC}
    \vspace{1.5cm}}}%
  }
  \caption{Organigramme de TAC-TIC — Groupe New Ways}
  \label{fig:organigramme}
\end{figure}

Notre stage s'est déroulé au sein du \textbf{département Développement informatique},
encadrés par le responsable technique de la division web/cloud.

% ============================================================
\section{Contexte et problématique}

\subsection{Contexte}

TAC-TIC gère simultanément une dizaine à une quinzaine de projets actifs, mobilisant
des équipes de 3 à 12 développeurs par projet. La coordination de ces équipes, la
traçabilité des décisions et le respect des délais constituent des défis quotidiens.
Les outils utilisés jusqu'alors — combinaison de Trello et de tableurs Excel — ne
permettaient pas :
\begin{itemize}
  \item une vue consolidée multi-projets ;
  \item la planification fine par sprints avec suivi de vélocité ;
  \item l'aide à la décision par l'IA pour l'assignation et le risque.
\end{itemize}

\subsection{Problématique}

\begin{mdframed}[linecolor=primary,linewidth=1.5pt,frametitle={}]
\centering
\textit{Comment concevoir une plateforme de gestion de projets qui intègre nativement
l'intelligence artificielle pour assister les chefs de projet dans leurs décisions
d'assignation, de planification et de gestion des risques, tout en garantissant une
expérience utilisateur fluide et une architecture logicielle maintenable ?}
\end{mdframed}

\subsection{Objectifs}

\textbf{Objectifs métier :}
\begin{itemize}
  \item Centraliser la gestion de tous les projets sur une seule plateforme ;
  \item Adopter les cérémonies Scrum (sprints, rétrospective, vélocité) ;
  \item Offrir des visualisations décisionnelles (Gantt, Kanban, dashboards) ;
  \item Automatiser les alertes de retard et les recommandations d'assignation.
\end{itemize}

\textbf{Objectifs techniques :}
\begin{itemize}
  \item Architecture distribuée (backend Laravel + SPA React + microservice ML) ;
  \item Sécurité multicouche (Sanctum, RBAC, policies, audit) ;
  \item Temps réel via WebSocket (Laravel Reverb) ;
  \item Qualité de code vérifiée par SonarCloud (Quality Gate \(\leq\) 3\,\% duplication).
\end{itemize}

% ============================================================
\section{Étude de l'existant et critique}

\subsection{Solutions analysées}

Nous avons étudié trois outils majeurs du marché : \textbf{Trello}, \textbf{Jira}
et \textbf{Asana}.

\begin{table}[H]
  \centering
  \caption{Comparaison des solutions de gestion de projets existantes}
  \label{tab:comparaison_existant}
  \small
  \begin{tabularx}{\textwidth}{L{4.5cm}C{2cm}C{2cm}C{2cm}C{2.5cm}}
    \toprule
    \thead{Fonctionnalité} & \thead{Trello} & \thead{Jira} & \thead{Asana} & \thead{TAC-TIC IntelliPM} \\
    \midrule
    \rowcolor{lightgray}
    Kanban drag-and-drop  & \cmark & \cmark & \cmark & \cmark \\
    Diagramme de Gantt    & \xmark & \cmark & \cmark & \cmark \\
    \rowcolor{lightgray}
    Chemin critique (CPM) & \xmark & \xmark & \xmark & \cmark \\
    Gestion des sprints   & \xmark & \cmark & \cmark & \cmark \\
    \rowcolor{lightgray}
    Recommandation ML IA  & \xmark & \xmark & \xmark & \cmark \\
    Estimation LLM        & \xmark & \xmark & \xmark & \cmark \\
    \rowcolor{lightgray}
    Notifications WebSocket & \xmark & \cmark & \cmark & \cmark \\
    Multi-rôles par projet & \cmark & \cmark & \cmark & \cmark \\
    \rowcolor{lightgray}
    Internationalisation FR & \cmark & \cmark & \cmark & \cmark \\
    Open source / hébergeable & \xmark & \xmark & \xmark & \cmark \\
    \bottomrule
  \end{tabularx}
\end{table}

\subsection{Critique}

\begin{itemize}
  \item \textbf{Trello} : simple et visuel, mais limité aux tableaux Kanban. Aucune
        notion de sprint, pas de Gantt, pas d'IA.
  \item \textbf{Jira} : puissant mais complexe à configurer. Coûteux pour les PME.
        Aucune recommandation ML native.
  \item \textbf{Asana} : bon équilibre fonctionnel, mais pas de chemin critique, pas
        d'IA intégrée, coût élevé par siège.
\end{itemize}

% ============================================================
\section{Solution proposée : TAC-TIC IntelliPM}

\begin{table}[H]
  \centering
  \caption{Apports différenciants de TAC-TIC IntelliPM}
  \label{tab:apports}
  \begin{tabularx}{\textwidth}{L{4.5cm}X}
    \toprule
    \thead{Fonctionnalité} & \thead{Apport} \\
    \midrule
    \rowcolor{lightgray}
    Recommandation ML       & Suggère le meilleur assigné pour chaque tâche (22 variables, AUC~0,92) \\
    Estimation LLM          & Durée estimée avec justification en langage naturel \\
    \rowcolor{lightgray}
    Détection de risques    & Score de risque tâche/projet + alertes proactives \\
    Chemin critique (CPM)   & Calcul automatique et visualisation sur le Gantt (barres rouges) \\
    \rowcolor{lightgray}
    Temps réel WebSocket    & Mises à jour Kanban et notifications sans rechargement \\
    RBAC à 2 niveaux        & Rôle système (admin/member/client) + rôle projet (chef/member/client) \\
    \rowcolor{lightgray}
    Journal d'audit         & Toutes les actions tracées avec IP, user, horodatage \\
    Dashboard multi-rôles   & Vue différenciée selon le profil (admin, member, client) \\
    \rowcolor{lightgray}
    Internationalisation    & Interface 100\,\% française, extensible \\
    Architecture ouverte    & Monorepo MIT, hébergeable en auto-managed \\
    \bottomrule
  \end{tabularx}
\end{table}

% ============================================================
\section{Méthode de travail : Scrum}

\subsection{Pourquoi Scrum ?}

La méthode Scrum a été retenue pour trois raisons principales :
\begin{enumerate}
  \item \textbf{Livraisons incrémentales} — chaque sprint livre un incrément fonctionnel
        démontrable, facilitant les retours du tuteur entreprise ;
  \item \textbf{Adaptabilité} — les backlogs sont révisés entre sprints pour intégrer
        les retours utilisateurs ;
  \item \textbf{Transparence} — les daily stand-ups et rétrospectives assurent une
        visibilité permanente sur l'avancement.
\end{enumerate}

\begin{figure}[H]
  \centering
  \IfFileExists{figures/cycle_scrum.png}{%
    \includegraphics[width=0.80\textwidth]{figures/cycle_scrum.png}%
  }{%
    \fbox{\parbox{0.80\textwidth}{\centering\vspace{1.5cm}
      \textbf{[IMAGE A INSÉRER]}\\[0.4em]
      \texttt{figures/cycle\_scrum.png}\\[0.3em]
      \textit{Cycle Scrum : Product Backlog → Sprint Backlog → Sprint → Incrément}
    \vspace{1.5cm}}}%
  }
  \caption{Cycle Scrum appliqué au projet TAC-TIC IntelliPM}
  \label{fig:cycle_scrum}
\end{figure}

\subsection{Planification des sprints}

\begin{table}[H]
  \centering
  \caption{Planification des 5 sprints du projet}
  \label{tab:sprints_planning}
  \begin{tabularx}{\textwidth}{C{1.5cm}L{4cm}X L{2.5cm}}
    \toprule
    \thead{Sprint} & \thead{Périmètre} & \thead{Livrables} & \thead{Durée} \\
    \midrule
    \rowcolor{lightgray}
    1 & Gestion des utilisateurs & Auth Sanctum, RBAC, profil, dashboard admin & 2 semaines \\
    2 & Projets et équipes       & CRUD projets, invitations, sprints, équipes & 2 semaines \\
    \rowcolor{lightgray}
    3 & Tâches et Kanban         & Wizard tâche, dépendances, Kanban, temps    & 2 semaines \\
    4 & Planification visuelle   & Gantt + CPM, calendrier, notifications      & 2 semaines \\
    \rowcolor{lightgray}
    5 & Intelligence artificielle & ML recommandation, LLM estimation, chat IA & 2 semaines \\
    \bottomrule
  \end{tabularx}
\end{table}

\section*{Conclusion}

Ce premier chapitre a délimité le cadre du projet : une entreprise dynamique, une
problématique réelle de coordination multi-projets, et une réponse technologique
ambitieuse articulant agilité, architecture distribuée et intelligence artificielle.
Le chapitre suivant entre dans le détail de l'étude préliminaire : acteurs, besoins,
product backlog et conception architecturale.
